\documentclass[aspectratio=169,11pt]{beamer}

% ------------------------------------------------------------
% Theme / style
% ------------------------------------------------------------
\usetheme{default}
\useinnertheme{rectangles}
\usefonttheme{professionalfonts}

\setbeamertemplate{navigation symbols}{}
\setbeamertemplate{footline}[frame number]

\usepackage{amsmath, amssymb, mathtools}
\usepackage{graphicx}
\usepackage{booktabs}
\usepackage{tikz}
\usepackage{xcolor}

% Optional simple color palette
\definecolor{darkgray}{HTML}{202124}
\definecolor{lightgray}{HTML}{F5F5F5}
\definecolor{accent}{HTML}{3B82F6}

\setbeamercolor{normal text}{fg=darkgray,bg=white}
\setbeamercolor{frametitle}{fg=darkgray,bg=white}
\setbeamercolor{title}{fg=darkgray}
\setbeamercolor{structure}{fg=accent}

\setbeamertemplate{itemize item}{\small$\blacktriangleright$}
\setbeamertemplate{itemize subitem}{\scriptsize$\bullet$}

% ------------------------------------------------------------
% Metadata
% ------------------------------------------------------------
\title{Operator-Aware World Model Compression}
\subtitle{Preserving planning behavior, not just prediction}
\author{Daniel Enriquez \and Man Sharma}
\institute{NYU Tandon}
\date{\today}

% ------------------------------------------------------------
% Convenience placeholders
% ------------------------------------------------------------
\newcommand{\placeholderfig}[2][0.75\linewidth]{%
    \begin{center}
    \fbox{
        \begin{minipage}[c][0.42\textheight][c]{#1}
            \centering
            \vspace{0.5em}
            \textit{#2}
            \vspace{0.5em}
        \end{minipage}
    }
    \end{center}
}

\newcommand{\Otheta}{\mathcal{O}_{\theta}}

% ------------------------------------------------------------
% Document
% ------------------------------------------------------------
\begin{document}

% ------------------------------------------------------------
\begin{frame}
    \titlepage
\end{frame}

% ------------------------------------------------------------
\begin{frame}{Motivation}
    % Goal: Explain why world-model compression matters.
    % Key message: These models are deployed inside planners, not used as passive predictors.

    \begin{itemize}
        \item World models are increasingly used for planning and control.
        \item Standard compression methods preserve weights, activations, or prediction error.
        \item But in planning, the model is used to choose actions.
    \end{itemize}

    % Optional figure: robot/world model/planner cartoon
    % \placeholderfig{World model inside CEM/MPC planner}
\end{frame}

% ------------------------------------------------------------
\begin{frame}{Problem}
    % Goal: State the mismatch between ordinary compression and planning use.

    \begin{block}{Standard compression question}
        Can we make the model smaller while preserving prediction accuracy?
    \end{block}

    \begin{block}{Our question}
        Can we make the model smaller while preserving the induced planning operator?
    \end{block}

    % Optional visual: prediction-MSE-preserved but action-ranking-changed diagram
\end{frame}

% ------------------------------------------------------------
\begin{frame}{Background: Latent World-Model Planning}
    % Goal: Explain the planning pipeline.

    \[
        h_t,\; g,\; A_{1:H}^{(1)}, \dots, A_{1:H}^{(N)}
        \quad \longrightarrow \quad
        \text{world model}
        \quad \longrightarrow \quad
        c_1,\dots,c_N
    \]

    \[
        A^* = \arg\min_i c_i
    \]

    % Optional figure: history + goal + candidate actions -> model -> costs -> CEM
    \placeholderfig{Planning pipeline diagram}
\end{frame}

% ------------------------------------------------------------
\begin{frame}{Models Studied}
    % Goal: Introduce LeWM and DINO-WM / SWM-PreJEPA.

    \begin{columns}[T,totalwidth=\textwidth]
        \begin{column}{0.48\textwidth}
            \textbf{LeWM}
            \begin{itemize}
                \item Latent world model from pixels.
                \item ViT-style encoder.
                \item Learned predictor / transition model.
                \item Used with CEM/MPC planning.
            \end{itemize}
        \end{column}

        \begin{column}{0.48\textwidth}
            \textbf{DINO-WM / SWM-PreJEPA}
            \begin{itemize}
                \item Frozen visual encoder.
                \item Learned Transformer predictor.
                \item Action/proprio conditioning.
                \item Same planning-style deployment.
            \end{itemize}
        \end{column}
    \end{columns}

    % Optional figure: side-by-side model architecture diagrams
\end{frame}

% ------------------------------------------------------------
\begin{frame}{Key Idea: Compress the Planning Operator}
    % Goal: Define the operator.

    \[
        \Otheta(h_t, g, A_{1:H}) = C_\theta(A_{1:H})
    \]

    \vspace{1em}

    \begin{itemize}
        \item The model induces a cost over candidate action sequences.
        \item CEM/MPC only needs the cost ranking to select actions.
        \item A compressed model may preserve planning even if latent predictions are imperfect.
    \end{itemize}
\end{frame}

% ------------------------------------------------------------
\begin{frame}{What Does Operator Preservation Mean?}
    % Goal: Define the metrics.

    \begin{itemize}
        \item \textbf{Rank correlation:} do teacher and student rank actions similarly?
        \item \textbf{Top-$k$ elite overlap:} do they agree on the CEM elite set?
        \item \textbf{Teacher regret:} how bad is the student's chosen action under the teacher?
        \item \textbf{First-action error:} does MPC choose a similar first action?
    \end{itemize}

    % Optional figure: teacher costs vs student costs / elite set overlap
\end{frame}

\begin{frame}{Why Prediction MSE Does Not Preserve Planning}
    \begin{block}{Planning is a ranking problem}
        The teacher selects
        \[
            A_T^* = \arg\min_A C_T(A).
        \]
        If the teacher action gap is
        \[
            \Delta =
            \min_{A \neq A_T^*}
            \left(C_T(A) - C_T(A_T^*)\right),
        \]
        then the student preserves the teacher action only if
        \[
            \sup_A |C_S(A)-C_T(A)| < \frac{\Delta}{2}.
        \]
    \end{block}

    \vspace{0.5em}

    \begin{itemize}
        \item Prediction MSE gives an average latent error, not a uniform cost-ranking guarantee.
        \item When two candidate actions have close costs, small prediction errors can flip the planner's choice.
        \item Therefore, compression should preserve action rankings and elite sets directly.
    \end{itemize}
\end{frame}

% ------------------------------------------------------------
\begin{frame}{Compression Setup}
    % Goal: Show teacher/student pipeline.

    \begin{columns}[T,totalwidth=\textwidth]
        \begin{column}{0.48\textwidth}
            \textbf{Teacher}
            \begin{itemize}
                \item Full LeWM or PreJEPA world model.
                \item Generates action costs and elite sets.
            \end{itemize}
        \end{column}

        \begin{column}{0.48\textwidth}
            \textbf{Student}
            \begin{itemize}
                \item Compressed predictor / transition operator.
                \item Encoder fixed initially.
                \item Trained or calibrated to preserve planning behavior.
            \end{itemize}
        \end{column}
    \end{columns}

    % Optional figure: teacher -> operator cache -> student
\end{frame}

% ------------------------------------------------------------
\begin{frame}{Compression Methods}
    % Goal: List baselines and our method.

    \begin{enumerate}
        \item Weight-SVD predictor compression.
        \item Activation-aware SVD.
        \item Prediction-only distillation.
        \item Operator-aware distillation.
    \end{enumerate}

    \vspace{1em}

    % Note to self: explain that SVD/pruning/quantization are standard tools;
    % the contribution is the operator-aware objective/evaluation.
\end{frame}

% ------------------------------------------------------------
\begin{frame}{Operator-Aware Loss}
    % Goal: Present the central loss.

    Teacher and student costs over the same candidate actions:
    \[
        c_T^{(i)} = \mathcal{O}_{\theta_T}(h_t,g,A_{1:H}^{(i)}),
        \qquad
        c_S^{(i)} = \mathcal{O}_{\theta_S}(h_t,g,A_{1:H}^{(i)}).
    \]

    Soft cost distribution:
    \[
        p_T(i) = \frac{\exp(-c_T^{(i)}/\tau)}
        {\sum_j \exp(-c_T^{(j)}/\tau)},
        \qquad
        p_S(i) = \frac{\exp(-c_S^{(i)}/\tau)}
        {\sum_j \exp(-c_S^{(j)}/\tau)}.
    \]

    \[
        \mathcal{L}_{\mathrm{op}}
        =
        D_{\mathrm{KL}}(p_T \,\|\, p_S)
        +
        \lambda \mathcal{L}_{\mathrm{elite}}.
    \]
\end{frame}

% ------------------------------------------------------------
\begin{frame}{Operator Cache}
    % Goal: Explain implementation object.

    \begin{itemize}
        \item Sample state/goal pairs from the offline dataset.
        \item Sample candidate action sequences.
        \item Store teacher costs, top-$k$ indices, elite sets, and selected actions.
        \item Reuse the same cache across compression methods.
    \end{itemize}

    % Optional figure/table: operator cache schema
    \placeholderfig{Operator cache schema}
\end{frame}

% ------------------------------------------------------------
\begin{frame}{Experimental Protocol}
    % Goal: State environments, models, and metrics.

    \begin{columns}[T,totalwidth=\textwidth]
        \begin{column}{0.33\textwidth}
            \textbf{Environments}
            \begin{itemize}
                \item TwoRoom
                \item PushT
                \item Cube
            \end{itemize}
        \end{column}

        \begin{column}{0.33\textwidth}
            \textbf{Models}
            \begin{itemize}
                \item LeWM
                \item SWM-PreJEPA
            \end{itemize}
        \end{column}

        \begin{column}{0.33\textwidth}
            \textbf{Metrics}
            \begin{itemize}
                \item Success rate
                \item Top-$k$ overlap
                \item Regret
                \item Model size
            \end{itemize}
        \end{column}
    \end{columns}
\end{frame}

% ------------------------------------------------------------
\begin{frame}{Main Result: Compression Frontier}
    % Goal: Placeholder for headline result.

    \placeholderfig{Compression ratio vs. planning success}

    % Expected interpretation:
    % Operator-aware compression should preserve planning success at higher compression ratios.
\end{frame}

% ------------------------------------------------------------
\begin{frame}{Diagnostics: Does Prediction Error Explain Planning?}
    % Goal: Show why operator metrics matter.

    \begin{columns}[T,totalwidth=\textwidth]
        \begin{column}{0.48\textwidth}
            \placeholderfig{Latent MSE vs. success}
        \end{column}

        \begin{column}{0.48\textwidth}
            \placeholderfig{Top-$k$ overlap / regret vs. success}
        \end{column}
    \end{columns}
\end{frame}

% ------------------------------------------------------------
\begin{frame}{Limitations}
    % Goal: Be honest.

    \begin{itemize}
        \item Operator cache depends on sampled states, goals, and action candidates.
        \item Preserving teacher behavior does not guarantee optimal behavior.
        \item Initial compression targets the predictor, not the visual encoder.
        \item Full benchmark cost depends on CEM settings and environment difficulty.
    \end{itemize}
\end{frame}

% ------------------------------------------------------------
\begin{frame}{Conclusion}
    % Goal: End with the claim.

    \begin{block}{Main idea}
        For world models used in planning, compression should preserve the induced action-ranking operator.
    \end{block}

    \begin{itemize}
        \item Standard compression asks whether the model still predicts.
        \item We ask whether the compressed model still chooses useful actions.
        \item Operator-aware compression targets rank, elite-set, and MPC-action preservation.
    \end{itemize}
\end{frame}

% ------------------------------------------------------------
\begin{frame}{Q \& A}
    \centering
    \Large Questions?
\end{frame}

% ------------------------------------------------------------
\end{document}